\chapter{Identity, Authentication, dan Authorization}
\label{ch:identitas}

\section*{Tujuan Pembelajaran}
\begin{itemize}
  \item Membedakan sesi yang disimpan server dari token bertanda tangan,
        lalu menyebut akibatnya pada pencabutan dan pada keadaan yang harus
        disimpan server.
  \item Merancang pasangan \textit{access token} dan \textit{refresh token}
        beserta rotasinya, lalu membuktikan pemakaian ulang penanda lama
        tertangkap.
  \item Menegakkan hak akses dengan RBAC dan ABAC sekaligus, lalu mengukur
        keputusannya sebagai matriks pengguna lawan alamat, termasuk isolasi
        antar-\textit{tenant}.
\end{itemize}

\section{Pendahuluan}

Bab 7 memasang pertahanan yang berlaku bagi semua orang: header, cookie yang
dibatasi, dan teks pengguna yang diperlakukan sebagai data. Bab ini
menanyakan hal yang berbeda, yaitu siapa yang sedang meminta, dan apa yang
boleh dilakukannya.

Dua pertanyaan itu sering disatukan padahal berbeda jauh.
\textit{Authentication} menjawab siapa, dan jawabannya satu kali di awal.
\textit{Authorization} menjawab boleh atau tidak, dan jawabannya diulang pada
tiap permintaan. Kesalahan paling mahal biasanya terjadi di pertanyaan kedua,
ketika aplikasi mengenali penggunanya dengan benar lalu lupa memeriksa apa
yang boleh dilihatnya.

Demonya kecil: satu aplikasi tanpa basis data, dengan empat pengguna pada dua
\textit{tenant}. Referensi utamanya adalah RFC 7519 untuk bentuk token
\cite{rfc7519}, RFC 6749 untuk kerangka OAuth 2.0 \cite{rfc6749}, RFC 9700
untuk praktik pengamanannya \cite{rfc9700}, dan NIST SP 800-207 untuk
gagasan \textit{zero trust} \cite{nist-zero-trust}.

\section{Cookie dan Session}

Sesi menyimpan keadaannya di server. Klien hanya membawa penanda acak di
dalam cookie, dan server menyimpan arti penanda itu di tempatnya sendiri.
Karena artinya ada di server, mencabut sesi cukup menghapus satu baris.

Analoginya seperti nomor antrean di loket. Kertasnya sendiri tidak memuat
apa-apa selain nomor, dan seluruh keterangan ada di buku petugas. Ketika
petugas mencoret satu nomor, kertas yang dipegang orangnya seketika tidak
berarti apa-apa lagi.

Demo bab ini menerbitkan sesi sepanjang 32 huruf, dengan atribut cookie yang
sama seperti Bab 7. Pencabutannya terbukti seketika, seperti pada Listing
\ref{lst:sesi10}.

\begin{lstlisting}[language=bash, caption={Sesi dicabut seketika saat keluar, tersedia di \berkas{sources/chapter10/aplikasi.py}}, label={lst:sesi10}]
curl -s -X POST "http://127.0.0.1:8070/masuk?cara=sesi&nama=budi&sandi=budi123" -D -
curl -s -H "Cookie: sesi=..." http://127.0.0.1:8070/api/pesanan
curl -s -X POST -H "Cookie: sesi=..." http://127.0.0.1:8070/keluar
curl -s -H "Cookie: sesi=..." http://127.0.0.1:8070/api/pesanan

# Keluarannya:
# Set-Cookie: sesi=4b2f8d7e...; HttpOnly; SameSite=Lax; Path=/
# {"tenant": "kampus-a", "jumlah": 2, "item": [...]}
# {"dicabut": true}
# {"pesan": "sesi tidak dikenal atau sudah dicabut"}
\end{lstlisting}

\section{JSON Web Token}

Token bertanda tangan menempuh jalan sebaliknya. Keterangannya dibawa di
dalam token itu sendiri, dan server tidak menyimpan apa pun. Keasliannya
dijaga tanda tangan, sehingga isinya boleh dibaca siapa saja tetapi tidak
dapat diubah tanpa ketahuan \cite{rfc7519}.

Listing \ref{lst:token10} memperlihatkan pemeriksaannya. Dua hal diperiksa
sekaligus, yaitu tanda tangan dan masa berlakunya.

\begin{lstlisting}[style=PythonStyle, caption={Memeriksa tanda tangan dan masa berlaku, tersedia di \berkas{sources/chapter10/token\_sederhana.py}}, label={lst:token10}]
def periksa(token):
  """Memeriksa tanda tangan dan masa berlaku, lalu mengembalikan isinya.

  Mengembalikan pasangan (isi, alasan). Isi bernilai None bila tokennya
  ditolak, dan alasannya menyebut penyebabnya, sehingga latihan dapat
  membedakan token palsu dari token yang sekadar kedaluwarsa.
  """
  bagian = token.split(".")
  if len(bagian) != 3:
    return None, "bentuk token salah"
  bagian_awal = f"{bagian[0]}.{bagian[1]}"
\end{lstlisting}

Dua akibat mengikuti pilihan itu. Pertama, server menjadi tidak berkeadaan,
sehingga permintaan dapat dilayani proses mana pun tanpa berbagi penyimpanan.
Kedua, token tidak dapat dicabut sebelum kedaluwarsa, karena tidak ada
tempat yang menyimpan daftar token yang masih berlaku.

Tabel \ref{tab:sesi-token} membandingkan keduanya dengan angka dari demo bab
ini.

\begin{table}[!htb]
  \centering
  \caption{Sesi lawan token bertanda tangan, hasil pengukuran nyata}
  \label{tab:sesi-token}
  \begin{tabularx}{\textwidth}{Xll}
    \hline
    \textbf{Hal} & \textbf{Sesi} & \textbf{Token} \\
    \hline
    Yang dibawa klien & penanda 32 huruf & token 169 byte \\
    \hline
    Yang disimpan server & arti tiap penanda & tidak ada \\
    \hline
    Dicabut saat keluar & seketika & tidak bisa, menunggu kedaluwarsa \\
    \hline
    Hilang saat proses mati & ya & tidak \\
    \hline
    Isinya terbaca klien & tidak & ya, tetapi tidak dapat diubah \\
    \hline
  \end{tabularx}
\end{table}

Baris terakhir sering disalahpahami. Token yang diubah satu huruf pun langsung
ditolak dengan alasan tanda tangan tidak cocok, tetapi isinya tetap terbaca
siapa pun yang memegangnya. Karena itu token tidak pernah diisi data rahasia.

\section{OAuth 2.x dan OpenID Connect}

OAuth 2.0 menjawab pertanyaan yang lebih luas daripada demo bab ini, yaitu
bagaimana satu aplikasi memperoleh izin mengakses data pengguna di aplikasi
lain, tanpa pengguna menyerahkan kata sandinya \cite{rfc6749}. Empat peran
terlibat di dalamnya, dan Gambar \ref{fig:oauth} merangkum alurnya.

\begin{figure}[!htb]
  \centering
  \begin{tikzpicture}[
    font=\scriptsize,
    kotak/.style={draw, rounded corners, minimum height=0.9cm,
                  minimum width=2.4cm, align=center, font=\scriptsize\bfseries},
    panah/.style={-{Stealth[length=2mm]}, thick}
  ]
    \node[kotak, fill=blue!8] (pengguna) at (0,0) {Pengguna\\{\tiny\mdseries resource owner}};
    \node[kotak, fill=orange!12] (klien) at (3.8,0) {Aplikasi\\{\tiny\mdseries client}};
    \node[kotak, fill=praditagreen!15] (otoritas) at (7.8,0.9) {Server otorisasi\\{\tiny\mdseries menerbitkan token}};
    \node[kotak, fill=green!10] (sumber) at (7.8,-0.9) {Server sumber\\{\tiny\mdseries memeriksa token}};

    \draw[panah] (pengguna) -- node[font=\tiny, above] {izin} (klien);
    \draw[panah] (klien.east) -- ++(0.4,0) |- (otoritas.west);
    \draw[panah] (klien.east) -- ++(0.4,0) |- (sumber.west);
    \node[font=\tiny, gray, anchor=east] at (7.3,0.72) {menukar kode dengan token};
    \node[font=\tiny, gray, anchor=east] at (7.3,-0.72) {memakai access token};
  \end{tikzpicture}
  \caption{Empat peran pada OAuth 2.0. Demo bab ini menyatukan server
           otorisasi dan server sumber dalam satu aplikasi}
  \label{fig:oauth}
\end{figure}

OpenID Connect menambahkan satu lapis di atas OAuth \cite{oidc-core}. OAuth
sendirian hanya menyatakan bahwa pemegang token boleh mengakses sesuatu,
tanpa menyebut siapa pemegangnya. OIDC menambahkan \textit{id token}, yaitu
token yang memang berisi identitas penggunanya.

Dua hal pada keduanya sering tertukar. Pembagiannya begini. OAuth menjawab
apa yang boleh diakses, sedangkan OIDC menjawab siapa yang sedang masuk.
Memakai \textit{access token} untuk menyimpulkan identitas adalah kesalahan
yang sering terjadi, dan RFC 9700 menyebutnya secara khusus
\cite{rfc9700}.

\section{Access Token, Refresh Token, dan Rotasi Token}

Dua token dipakai bersama karena keduanya menanggung risiko yang berbeda.
\textit{Access token} ikut terkirim pada tiap permintaan, sehingga umurnya
dibuat pendek. \textit{Refresh token} hanya dikirim ke satu alamat, sehingga
umurnya boleh panjang.

Tabel \ref{tab:umur-token} memuat perilaku \textit{access token} pada demo
bab ini, yang umurnya sengaja hanya lima detik.

\begin{table}[!htb]
  \centering
  \caption{Umur \textit{access token} lima detik, hasil pengukuran nyata}
  \label{tab:umur-token}
  \begin{tabularx}{\textwidth}{Xll}
    \hline
    \textbf{Kapan dipakai} & \textbf{Status} & \textbf{Alasan} \\
    \hline
    Segera sesudah terbit & 200 & sah \\
    \hline
    Sesudah 2 detik & 200 & sah \\
    \hline
    Sesudah 6 detik & 401 & token kedaluwarsa \\
    \hline
    Diubah satu huruf & 401 & tanda tangan tidak cocok \\
    \hline
  \end{tabularx}
\end{table}

Rotasi menjawab satu kelemahan \textit{refresh token} yang berumur panjang:
bila penanda itu bocor, pencurinya dapat memakainya berkali-kali. Dengan
rotasi, tiap penyegaran menerbitkan penanda baru dan menandai yang lama
terpakai. Bila yang lama dipakai lagi, keadaan itu hanya mungkin terjadi
karena ada dua pihak yang memegangnya \cite{rfc9700}.

Listing \ref{lst:rotasi} memperlihatkan jawabannya, yaitu mencabut seluruh
keluarga penanda tersebut.

\begin{lstlisting}[style=PythonStyle, caption={Pemakaian ulang menandakan penanda sudah di tangan lain, tersedia di \berkas{sources/chapter10/aplikasi.py}}, label={lst:rotasi}]
def cabut_keluarga(keluarga):
  """Mencabut seluruh refresh token satu keluarga sekaligus.

  Dipanggil ketika ada refresh token lama yang dipakai ulang. Keadaan itu
  menandakan penanda tersebut sudah berada di tangan lain, dan cara teraman
  adalah menghentikan seluruh rantainya.
  """
  jumlah = 0
  for penanda, catatan in list(REFRESH.items()):
    if catatan["keluarga"] == keluarga:
      del REFRESH[penanda]
      jumlah += 1
  return jumlah
\end{lstlisting}

Tabel \ref{tab:rotasi} memuat ketiga langkah tersebut, diukur berurutan pada
demo bab ini.

\begin{table}[!htb]
  \centering
  \caption{Rotasi \textit{refresh token} dan penangkapan pemakaian ulang,
           hasil pengukuran nyata}
  \label{tab:rotasi}
  \begin{tabularx}{\textwidth}{Xll}
    \hline
    \textbf{Langkah} & \textbf{Status} & \textbf{Hasil} \\
    \hline
    Menyegarkan dengan penanda yang sah & 200 & terbit pasangan token baru \\
    \hline
    Memakai ulang penanda lama & 401 & dipakai ulang, 2 penanda dicabut \\
    \hline
    Memakai penanda baru sesudah pencabutan & 401 & penanda tidak dikenal \\
    \hline
  \end{tabularx}
\end{table}

Baris ketiga itu yang sering mengejutkan. Penanda baru ikut mati karena
seluruh keluarganya dicabut, dan pengguna yang sah harus masuk ulang. Harga
itu memang disengaja: ketika satu penanda terbukti bocor, tidak ada cara
membedakan pemilik aslinya dari pemegang lainnya.

\section{RBAC dan ABAC}

RBAC memutuskan dari peran. Tiap peran membawa daftar izin, dan yang
diperiksa adalah izinnya, bukan nama perannya. ABAC memutuskan dari sifat
permintaan: siapa pemohonnya, siapa pemilik datanya, dan dari
\textit{tenant} mana keduanya.

Listing \ref{lst:abac} memperlihatkan keduanya bekerja dalam satu fungsi.
Urutannya penting: pemeriksaan \textit{tenant} berada paling depan, sehingga
peran setinggi apa pun tidak dapat menembus batas itu.

\begin{lstlisting}[style=PythonStyle, caption={Tiga sifat diperiksa berurutan, tersedia di \berkas{sources/chapter10/kebijakan.py}}, label={lst:abac}]
def boleh_baca_pesanan(identitas, pesanan):
  """Memutuskan satu pesanan boleh dibaca atau tidak, secara ABAC.

  Tiga sifat ikut dibaca: tenant pemohon, tenant pesanan, dan pemiliknya.
  Mengembalikan pasangan (boleh, alasan), sehingga alasan penolakan dapat
  ikut diuji dan ikut dicatat.
  """
  if identitas["tenant"] != pesanan["tenant"]:
    return False, "tenant berbeda"
  if punya_izin(identitas["peran"], "baca_semua_pesanan"):
    return True, "peran boleh membaca seluruh pesanan tenantnya"
  if pesanan["pemilik"] == identitas["sub"]:
    return True, "pemohon adalah pemilik pesanan"
  return False, "bukan pemilik dan perannya tidak mencakup"
\end{lstlisting}

\section{Multi-Tenancy dan Isolasi Data}

\textit{Multi-tenancy} berarti satu aplikasi melayani beberapa penyewa
sekaligus, dan tiap penyewa hanya boleh melihat datanya sendiri. Pada demo
bab ini ada dua \textit{tenant}, yaitu kampus-a dan kampus-b, dan tiap
pesanan membawa nama \textit{tenant}-nya.

Tabel \ref{tab:matriks} memuat keputusan seluruh persilangan empat pengguna
dengan empat alamat, diukur sekali jalan.

\begin{table}[!htb]
  \centering
  \caption{Matriks hak akses empat pengguna lawan empat alamat, hasil
           pengukuran nyata}
  \label{tab:matriks}
  \begin{tabularx}{\textwidth}{Xllll}
    \hline
    \textbf{Pengguna} & \textbf{/api/pesanan} & \textbf{/pesanan/1} & \textbf{/pesanan/3} & \textbf{/laporan} \\
    \hline
    ani, admin kampus-a & 200, 2 baris & 200 & 404 & 200 \\
    \hline
    budi, staf kampus-a & 200, 2 baris & 200 & 404 & 403 \\
    \hline
    cita, pembeli kampus-a & 200, 1 baris & 200 & 404 & 403 \\
    \hline
    dedi, staf kampus-b & 200, 1 baris & 404 & 200 & 403 \\
    \hline
  \end{tabularx}
\end{table}

Empat belas sel itu menjelaskan seluruh bab ini. Kolom terakhir murni RBAC:
hanya ani yang punya izin \texttt{baca\_laporan}, sehingga tiga lainnya
ditolak 403. Kolom pertama murni ABAC: cita hanya melihat satu baris, yaitu
pesanan miliknya, sedangkan budi dan ani melihat dua baris milik
\textit{tenant} mereka.

Dua kolom tengah memuat isolasi \textit{tenant}. Pesanan nomor 1 milik
kampus-a dijawab 200 bagi ketiga penggunanya dan 404 bagi dedi, sedangkan
pesanan nomor 3 milik kampus-b berlaku sebaliknya.

Pilihan status itu disengaja dan layak diperhatikan. Permintaan lintas
\textit{tenant} dijawab 404, bukan 403, karena 403 berarti "ada, tetapi Anda
tidak boleh", dan jawaban itu sendiri sudah membocorkan keberadaan datanya.
Penolakan di dalam satu \textit{tenant} tetap 403, karena di sana keberadaan
pesanan memang bukan rahasia.

\section{Identity Provider dan Konsep Zero Trust}

\textit{Identity provider} adalah satu tempat yang memegang identitas
pengguna untuk banyak aplikasi sekaligus. Dengan begitu, kata sandi hanya
diketahui satu sistem, dan aplikasi lain cukup mempercayai token yang
diterbitkannya.

\textit{Zero trust} adalah cara memandang jaringan: tidak ada permintaan yang
dipercaya hanya karena datang dari dalam \cite{nist-zero-trust}. Setiap
permintaan membawa identitasnya sendiri dan diperiksa ulang di tempat
tujuannya.

Tabel \ref{tab:zero-trust} membandingkan dua cara pandang tersebut.

\begin{table}[!htb]
  \centering
  \caption{Cara pandang lama lawan \textit{zero trust}}
  \label{tab:zero-trust}
  \begin{tabularx}{\textwidth}{lXX}
    \hline
    \textbf{Hal} & \textbf{Cara lama} & \textbf{Zero trust} \\
    \hline
    Dasar kepercayaan & Letak di jaringan & Identitas tiap permintaan \\
    \hline
    Pemeriksaan & Di pintu masuk & Di tiap layanan \\
    \hline
    Permintaan dari dalam & Dipercaya & Diperiksa sama ketatnya \\
    \hline
    Akibat satu titik jebol & Seluruh jaringan terbuka & Terbatas pada hak penanda itu \\
    \hline
  \end{tabularx}
\end{table}

Demo bab ini mengikuti cara kedua, walaupun kecil. Tidak ada alamat yang
melayani permintaan tanpa penanda, dan pemeriksaan hak aksesnya terjadi di
tiap alamat, bukan sekali di depan.

\section{Ringkasan}

Sesi dan token menjawab pertanyaan yang sama dengan harga yang berbeda. Sesi
menyimpan artinya di server, sehingga pencabutannya seketika dan terbukti
pada demo ini: sesudah keluar, penanda yang sama langsung ditolak. Token
tidak menyimpan apa pun di server, sehingga 169 byte yang dibawa klien sudah
cukup, tetapi pencabutannya harus menunggu kedaluwarsa.

Umur pendek menutup sebagian kelemahan itu. \textit{Access token} pada demo
ini berlaku lima detik: sah segera sesudah terbit, sah pada detik kedua,
lalu ditolak pada detik keenam. Token yang diubah satu huruf ditolak lebih
awal lagi, karena tanda tangannya tidak cocok.

Rotasi menutup sisanya. Tiap penyegaran menerbitkan penanda baru dan
menandai yang lama terpakai, dan ketika penanda lama dipakai ulang, seluruh
keluarganya dicabut. Pada demo ini dua penanda hilang sekaligus, termasuk
penanda baru yang sah, dan pemiliknya harus masuk ulang.

Hak akses diputuskan dua model yang bekerja bersama. RBAC menjawab dari izin
yang dibawa peran, dan hanya satu dari empat pengguna yang boleh membuka
laporan. ABAC menjawab dari sifat permintaan, sehingga pembeli hanya melihat
satu pesanan miliknya sendiri sedangkan staf melihat dua.

Batas \textit{tenant} diperiksa paling depan, sebelum peran dan sebelum
kepemilikan. Permintaan lintas \textit{tenant} dijawab 404, bukan 403, karena
jawaban 403 sudah menyatakan bahwa datanya ada. Dalam sistem yang melayani
banyak penyewa, keberadaan data pun termasuk yang dijaga.

\section{Latihan}

Seluruh file yang dipakai pada latihan ini tersedia di
\berkas{sources/chapter10/}, dihitung relatif terhadap akar repositori.
Latihan ini hanya membutuhkan Python 3, tanpa Docker dan tanpa basis data.

Seluruh permintaan diarahkan ke aplikasi yang berjalan di komputer sendiri.
Kata sandi pada \berkas{sources/chapter10/pengguna.json} hanya untuk latihan,
dan kunci tandanya dibaca dari \textit{environment} \texttt{KUNCI\_TOKEN}.
Hasil tiap mahasiswa akan sama pada bagian status, dan berbeda pada bagian
penanda yang memang acak.

Penyiapannya satu perintah, seperti pada Listing \ref{lst:c10-siapkan}.

\begin{lstlisting}[language=bash, caption={Menyalakan aplikasi identitas, tersedia di \berkas{sources/chapter10/aplikasi.py}}, label={lst:c10-siapkan}]
cd sources/chapter10
KUNCI_TOKEN=kunci-latihan python3 aplikasi.py &

# Keluarannya:
# Melayani di http://127.0.0.1:8070/  (Ctrl+C untuk berhenti)
\end{lstlisting}

\subsection*{Latihan 1: Sesi lawan Token}

Latihan ini membuktikan bahwa pencabutan seketika hanya dimiliki sesi, dan
bahwa isi token terbaca siapa pun yang memegangnya.

\begin{enumerate}
  \item Masuk dengan cara sesi, catat panjang penandanya dan atribut
        cookienya, seperti pada Listing \ref{lst:c10-lat1-jalan}.
  \item Pakai sesi itu untuk membuka daftar pesanan, lalu keluar, lalu
        pakai sesi yang sama sekali lagi. Catat ketiga hasilnya.
  \item Masuk dengan cara token, lalu dekode bagian tengah tokennya dengan
        \texttt{base64}. Catat apa saja yang terbaca di dalamnya.
  \item Isi Tabel \ref{tab:c10-lat1-hasil}. Tabel \ref{tab:c10-lat1-contoh}
        memperlihatkan contoh yang sudah terisi.
  \item Jawab tiga pertanyaan berikut. Pertama, mengapa token tidak dapat
        dicabut seketika? Kedua, apa yang hilang bila seluruh keterangan
        pengguna dimasukkan ke dalam token? Ketiga, kapan sesi lebih pantas
        dipakai daripada token?
\end{enumerate}

\begin{lstlisting}[language=bash, caption={Membandingkan sesi dan token, tersedia di \berkas{sources/chapter10/aplikasi.py}}, label={lst:c10-lat1-jalan}]
curl -s -D - -X POST "http://127.0.0.1:8070/masuk?cara=sesi&nama=budi&sandi=budi123"
curl -s -X POST "http://127.0.0.1:8070/masuk?cara=token&nama=budi&sandi=budi123"

# Keluarannya:
# Set-Cookie: sesi=4b2f8d7e...; HttpOnly; SameSite=Lax; Path=/
# {"cara": "sesi", "panjang_penanda": 32}
# {"cara": "token", "access_token": "eyJhbGciOi...", "panjang_access": 169}
\end{lstlisting}

\begin{table}[!htb]
  \centering
  \caption{Contoh hasil Latihan 1 yang sudah terisi}
  \label{tab:c10-lat1-contoh}
  \begin{tabularx}{\textwidth}{Xll}
    \hline
    \textbf{Hal} & \textbf{Sesi} & \textbf{Token} \\
    \hline
    Panjang yang dibawa klien & 32 huruf & 169 byte \\
    \hline
    Sesudah keluar & ditolak & masih berlaku \\
    \hline
  \end{tabularx}
\end{table}

\begin{table}[!htb]
  \centering
  \caption{Lembar isian hasil pengamatan Latihan 1}
  \label{tab:c10-lat1-hasil}
  \begin{tabularx}{\textwidth}{Xll}
    \hline
    \textbf{Hal} & \textbf{Sesi} & \textbf{Token} \\
    \hline
    Panjang yang dibawa klien & \dots & \dots \\
    \hline
    Sesudah keluar & \dots & \dots \\
    \hline
    Isi yang terbaca klien & \dots & \dots \\
    \hline
    Yang disimpan server & \dots & \dots \\
    \hline
  \end{tabularx}
\end{table}

\subsection*{Latihan 2: Umur dan Rotasi Token}

Latihan ini membuktikan bahwa umur pendek dan rotasi menjawab dua risiko yang
berbeda.

\begin{enumerate}
  \item Masuk dengan cara token, lalu pakai \textit{access token}-nya
        segera, pada detik kedua, dan pada detik keenam. Catat ketiga
        statusnya.
  \item Ubah satu huruf terakhir token itu, lalu pakai lagi dan catat
        alasan penolakannya.
  \item Segarkan dengan \textit{refresh token}, catat pasangan baru yang
        terbit, lalu pakai ulang penanda lama seperti pada Listing
        \ref{lst:c10-lat2-jalan}.
  \item Coba penanda baru sesudah pemakaian ulang itu, lalu isi Tabel
        \ref{tab:c10-lat2-hasil}. Tabel \ref{tab:c10-lat2-contoh}
        memperlihatkan contoh yang sudah terisi.
  \item Jawab dua pertanyaan berikut. Pertama, mengapa penanda baru ikut
        mati? Kedua, apa yang terjadi bila rotasi dipakai tanpa penangkapan
        pemakaian ulang?
\end{enumerate}

\begin{lstlisting}[language=bash, caption={Rotasi dan pemakaian ulang refresh token, tersedia di \berkas{sources/chapter10/aplikasi.py}}, label={lst:c10-lat2-jalan}]
curl -s -X POST "http://127.0.0.1:8070/segarkan?refresh=PENANDA_LAMA"
curl -s -X POST "http://127.0.0.1:8070/segarkan?refresh=PENANDA_LAMA"
curl -s -X POST "http://127.0.0.1:8070/segarkan?refresh=PENANDA_BARU"

# Keluarannya:
# {"access_token": "eyJhbGciOi...", "refresh_token": "50dbfff1..."}
# {"pesan": "refresh token dipakai ulang", "dicabut": 2}
# {"pesan": "refresh token tidak dikenal"}
\end{lstlisting}

\begin{table}[!htb]
  \centering
  \caption{Contoh hasil Latihan 2 yang sudah terisi}
  \label{tab:c10-lat2-contoh}
  \begin{tabularx}{\textwidth}{Xll}
    \hline
    \textbf{Langkah} & \textbf{Status} & \textbf{Alasan} \\
    \hline
    Token sesudah 6 detik & 401 & token kedaluwarsa \\
    \hline
    Memakai ulang penanda lama & 401 & dipakai ulang, 2 dicabut \\
    \hline
  \end{tabularx}
\end{table}

\begin{table}[!htb]
  \centering
  \caption{Lembar isian hasil pengamatan Latihan 2}
  \label{tab:c10-lat2-hasil}
  \begin{tabularx}{\textwidth}{Xll}
    \hline
    \textbf{Langkah} & \textbf{Status} & \textbf{Alasan} \\
    \hline
    Token segera sesudah terbit & \dots & \dots \\
    \hline
    Token sesudah 2 detik & \dots & \dots \\
    \hline
    Token sesudah 6 detik & \dots & \dots \\
    \hline
    Token yang diubah satu huruf & \dots & \dots \\
    \hline
    Memakai ulang penanda lama & \dots & \dots \\
    \hline
    Penanda baru sesudah pencabutan & \dots & \dots \\
    \hline
  \end{tabularx}
\end{table}

\subsection*{Latihan 3: Matriks Hak Akses}

Latihan ini membuktikan bahwa RBAC, ABAC, dan isolasi \textit{tenant}
terbaca sekaligus dari satu tabel.

\begin{enumerate}
  \item Jalankan alat ukurnya seperti pada Listing \ref{lst:c10-lat3-jalan},
        lalu salin hasilnya ke Tabel \ref{tab:c10-lat3-hasil}. Tabel
        \ref{tab:c10-lat3-contoh} memperlihatkan contoh yang sudah terisi.
  \item Tandai sel mana yang ditentukan peran, sel mana yang ditentukan
        kepemilikan, dan sel mana yang ditentukan \textit{tenant}.
  \item Tambahkan satu izin \texttt{baca\_laporan} pada peran staf di
        \berkas{sources/chapter10/kebijakan.py}, jalankan ulang, lalu catat
        sel mana yang berubah.
  \item Jawab tiga pertanyaan berikut. Pertama, mengapa permintaan lintas
        \textit{tenant} dijawab 404 dan bukan 403? Kedua, mengapa
        pemeriksaan \textit{tenant} ditaruh paling depan? Ketiga, apa yang
        terjadi bila penyaringan daftar dan pemeriksaan rincian memakai
        aturan yang berbeda?
\end{enumerate}

\begin{lstlisting}[language=bash, caption={Menjalankan matriks hak akses, tersedia di \berkas{sources/chapter10/ukur-akses.py}}, label={lst:c10-lat3-jalan}]
python3 ukur-akses.py

# Keluarannya:
#   Pengguna    /api/pesanan  /api/pesanan/1  /api/pesanan/3    /api/laporan
#   ani         200, 2 baris             200             404             200
#   budi        200, 2 baris             200             404             403
#   cita        200, 1 baris             200             404             403
#   dedi        200, 1 baris             404             200             403
\end{lstlisting}

\begin{table}[!htb]
  \centering
  \caption{Contoh hasil Latihan 3 yang sudah terisi}
  \label{tab:c10-lat3-contoh}
  \begin{tabularx}{\textwidth}{Xllll}
    \hline
    \textbf{Pengguna} & \textbf{Daftar} & \textbf{/1} & \textbf{/3} & \textbf{Laporan} \\
    \hline
    ani, admin kampus-a & 200, 2 baris & 200 & 404 & 200 \\
    \hline
    dedi, staf kampus-b & 200, 1 baris & 404 & 200 & 403 \\
    \hline
  \end{tabularx}
\end{table}

\begin{table}[!htb]
  \centering
  \caption{Lembar isian hasil pengukuran Latihan 3}
  \label{tab:c10-lat3-hasil}
  \begin{tabularx}{\textwidth}{Xllll}
    \hline
    \textbf{Pengguna} & \textbf{Daftar} & \textbf{/1} & \textbf{/3} & \textbf{Laporan} \\
    \hline
    ani, admin kampus-a & \dots & \dots & \dots & \dots \\
    \hline
    budi, staf kampus-a & \dots & \dots & \dots & \dots \\
    \hline
    cita, pembeli kampus-a & \dots & \dots & \dots & \dots \\
    \hline
    dedi, staf kampus-b & \dots & \dots & \dots & \dots \\
    \hline
    Sesudah staf diberi izin laporan & \dots & \dots & \dots & \dots \\
    \hline
  \end{tabularx}
\end{table}

\subsection*{Latihan 4: Menambah Tenant dan Peran Baru}

Latihan ini menguji apakah aturan hak akses benar-benar terpusat, atau
sebenarnya tersebar di banyak tempat.

\begin{enumerate}
  \item Tambahkan satu \textit{tenant} ketiga beserta satu penggunanya dan
        satu pesanannya ke
        \berkas{sources/chapter10/pengguna.json}.
  \item Tambahkan peran baru, misalnya auditor yang boleh membaca laporan
        tetapi tidak boleh membaca pesanan, ke
        \berkas{sources/chapter10/kebijakan.py}.
  \item Tambahkan kedua pengguna baru itu ke daftar pada
        \berkas{sources/chapter10/ukur-akses.py}, jalankan ulang, lalu
        catat matriksnya ke Tabel \ref{tab:c10-lat4-hasil}. Tabel
        \ref{tab:c10-lat4-contoh} memperlihatkan bentuk isian yang
        diharapkan.
  \item Catat berapa file yang harus disunting untuk kedua penambahan itu.
  \item Jawab dua pertanyaan berikut. Pertama, berapa file yang berubah, dan
        apakah jumlah itu menandakan aturannya terpusat? Kedua, apa yang
        perlu diubah bila aturannya berganti dari peran menjadi izin per
        pengguna?
\end{enumerate}

\begin{table}[!htb]
  \centering
  \caption{Bentuk isian Latihan 4 yang diharapkan}
  \label{tab:c10-lat4-contoh}
  \begin{tabularx}{\textwidth}{Xlll}
    \hline
    \textbf{Pengguna baru} & \textbf{Daftar} & \textbf{Rincian} & \textbf{Laporan} \\
    \hline
    auditor kampus-a & 200, 0 baris & 403 & 200 \\
    \hline
  \end{tabularx}
\end{table}

\begin{table}[!htb]
  \centering
  \caption{Lembar isian hasil pengamatan Latihan 4}
  \label{tab:c10-lat4-hasil}
  \begin{tabularx}{\textwidth}{Xlll}
    \hline
    \textbf{Pengguna baru} & \textbf{Daftar} & \textbf{Rincian} & \textbf{Laporan} \\
    \hline
    Pengguna \textit{tenant} ketiga & \dots & \dots & \dots \\
    \hline
    Peran auditor & \dots & \dots & \dots \\
    \hline
    Jumlah file yang disunting & \multicolumn{3}{l}{\dots} \\
    \hline
  \end{tabularx}
\end{table}

Setelah seluruh latihan selesai, aplikasinya dihentikan dengan Ctrl+C. Tidak
ada container yang perlu dimatikan, karena bab ini memang tidak memakainya.
